%! Author = layne
%! Date = 4/12/26

% Preamble
%  preamble
\documentclass[12pt]{article}

% dependencies
\usepackage[margin=0.6in]{geometry}
\usepackage{xcolor}
\usepackage[most]{tcolorbox}
\usepackage{enumitem}
\usepackage{multicol}
\usepackage{amsmath,amssymb}
\usepackage{array}
\usepackage{tabularx}
\usepackage{unicode-math}
\usepackage{graphicx}
\usepackage{pdfpages}
\usepackage{ltablex}
\usepackage{lettrine}
\keepXColumns
\tcbuselibrary{breakable}

% definitions
\definecolor{sky}{HTML}{EAF3FB}
\definecolor{navy}{HTML}{1F3A5F}
\definecolor{redbox}{HTML}{FCECEC}
\definecolor{greenbox}{HTML}{EDF8F0}
\definecolor{goldbox}{HTML}{FFF7E6}
\newtcolorbox{skillbox}[2][]{
    colback=#2,
    colframe=navy,
    boxrule=0.8pt,
    arc=2mm,
    left=2mm,right=2mm,top=1.5mm,bottom=1.5mm,
    fonttitle=\bfseries,
    title={#1},
    halign title=center,
    breakable
}
\newtcolorbox{fixedskillbox}[2][]{
    colback=#2,
    colframe=navy,
    boxrule=0.8pt,
    arc=2mm,
    left=2mm,right=2mm,top=0.3mm,bottom=1mm,
    fonttitle=\bfseries,
    title={#1},
    halign title=center,
    unbreakable
}

\setlist[itemize]{leftmargin=1.2em,itemsep=0.35em,topsep=0.35em}
\setlist[enumerate]{leftmargin=1.4em,itemsep=0.35em,topsep=0.35em}

% Centered column types
\newcolumntype{C}[1]{>{\centering\arraybackslash}p{#1}}
\newcolumntype{Y}{>{\centering\arraybackslash}X}

\newcommand{\CourseName}{Algebra 2: Shepherd}
\newcommand{\SchoolYear}{2026--2027}
\newcommand{\LoremIpsum}{
    Lorem ipsum dolor sit amet, consectetur adipiscing elit, sed do eiusmod
    tempor incididunt ut labore et dolore magna aliqua. Ut enim ad minim
    veniam, quis nostrud exercitation ullamco laboris nisi ut aliquip ex ea
    commodo consequat. Duis aute irure dolor in reprehenderit in voluptate
    velit esse cillum dolore eu fugiat nulla pariatur. Excepteur sint
    occaecat cupidatat non proident, sunt in culpa qui officia deserunt
    mollit anim id est laborum.
}
\newcommand{\TallMath}[1]{$\displaystyle #1\rule[-1.4em]{0pt}{3.2em}$}

\begin{document}
    \begin{center}
        {\Large\bfseries \CourseName \quad \SchoolYear}
    \end{center}

    \renewcommand{\arraystretch}{1.6}
    \begin{skillbox}[Welcome to Algebra 2!]{sky}
        \setlength{\parskip}{0.6em}
        \setlength{\parindent}{0pt}
        \setlength{\columnsep}{1cm}
        \begin{multicols}{2}
            \raggedcolumns
            \lettrine{I} am excited to welcome you to class and begin what I hope will be
            a challenging, rewarding, and successful year together. Algebra 2 is
            a course where we will build powerful problem-solving skills, strengthen
            logical thinking, and explore the many ways mathematics helps us
            understand patterns and relationships in the world around us.

            \lettrine{F}{or} some of you, math may already be a favorite subject. For others,
            math may feel stressful, frustrating, or intimidating. If you think
            you “aren’t a math person” or you have had difficult experiences
            with math in the past, I want to ask you to leave those experiences at the door
            and approach this class with an open mind--a beginner's mind.  And we are ALL
            beginners!
            \vspace{1em}
            \begin{center}
            \itshape
            “In the beginner’s mind there are many possibilities,
            but in the expert’s there are few.”
            \vspace{0.5em}
            --- Shunryu Suzuki, \textit{Zen Mind, Beginner’s Mind}
            \end{center}

            \vspace{1em}
            \lettrine{T}{his} classroom is a place where:
            \begin{itemize}
                \item mistakes are part of learning,
                \item questions are encouraged,
                \item effort matters,
                \item and growth is expected.
            \end{itemize}

            \lettrine{S}{uccess} in Algebra 2 is not about being naturally gifted at math.
            It is about persistence, curiosity, practice, and a willingness
            to keep trying when something feels difficult.

            \lettrine{T}{hroughout} the year, we will focus not only on finding answers,
            but also on understanding patterns, making connections,
            communicating ideas, and developing confidence in your own thinking.
            Some days will be challenging, but you will not be expected
            to figure everything out alone. We will work together as a class,
            ask good questions, learn from mistakes, and celebrate progress
            along the way.

            \lettrine{O}{n} the first day of class, we will go over important course information
            including:
            \begin{itemize}
                \item the syllabus,
                \item classroom expectations,
                \item materials and routines,
                \item and what you can expect from me as your teacher.
            \end{itemize}

            \lettrine{M}{ost} importantly, we will begin building a classroom community
            where everyone feels supported and capable of learning mathematics.

            \lettrine{I} am genuinely looking forward to getting to know each of you
            and learning with you this year.
        \end{multicols}
    \end{skillbox}

    \renewcommand{\arraystretch}{1.6}
    \begin{skillbox}[Important information]{goldbox}
        \renewcommand{\arraystretch}{1.35}
        \begin{tabularx}{\linewidth}{
                >{\bfseries}p{0.24\linewidth}
            X
        }
        \endhead
        {Contacting me}
        & {You may contact me at any time at my email address, however,
        I need to set some clear expectations around response times:
        I will get back to you within 48 hours.
        I will most likely get back to you before that time, but there may be
        weeks where I need that extra time to get together a meaningful response to your question(s).}
        \\[1em]
        \hline

        {Grading scale}
        &
        \begin{minipage}[t]{\linewidth}
            \begin{itemize}
                \item \textbf{A}: $90-100$
                \item \textbf{B}: $80-89$
                \item \textbf{C}: $70-79$
                \item \textbf{D}: $60-69$
                \item \textbf{F}: $0-59$
            \end{itemize}
        I will update grades at least once per week, usually on weekends.
        Your assignments will be weighted as follows:
        \begin{itemize}
            \item \textbf{Bronze}: Homework, Warm-ups, Exit Tickets -- 20\% of grade
            \item \textbf{Silver}: Quizzes, Assignments -- 30\% of grade
            \item \textbf{Gold}: Tests, Exams, Projects -- 50\% of grade
        \end{itemize}
        \end{minipage}
        \\ \hline
        \end{tabularx}
    \end{skillbox}
\end{document}